\documentclass[11pt,a4paper]{article}

\usepackage[margin=1in]{geometry}
\usepackage{amsmath,amssymb}
\usepackage{booktabs}
\usepackage{hyperref}
\usepackage{xcolor}

\title{A Note on the Z-Plane Covariance Structure\\
of the Hilbert Space-Filling Curve}

\author{Rick Wesson\\
Support Intelligence, Inc.}
\date{June 2026}

\begin{document}
\maketitle

\begin{abstract}
We observe that the Z-plane spatial covariance matrix of the
$D$-dimensional Hilbert curve is tridiagonal: all entries $C[z, z \pm d]$
for $d > 1$ are identically zero. This follows directly from
face-adjacency --- no spatial neighbor pair connects Z-planes separated by
more than one unit. Because the matrix is also Toeplitz (constant diagonal
and off-diagonal), its eigendecomposition is the known result for
tridiagonal Toeplitz matrices \cite{grenander1958}: $\lambda_k = c_0 +
2c_1\cos(\pi k/(n+1))$ with eigenvectors $v_k[j] = \sin(\pi k j/(n+1))$.
The formula matches numerical computation at $r > 0.98$ across orders
4--7. Practical consequences include a diagnostic for spatial continuity
of space-filling curves and an exact bound on cross-plane correlation for
spatial indexing.
\end{abstract}

% ============================================================================
\section{The Observation}
% ============================================================================

The Hilbert curve $H_D: \mathbb{N} \to \mathbb{Z}^D$ \cite{hilbert1891}
maps consecutive integers to face-adjacent cells in a $D$-dimensional
hypercube of side $2^n$. Project onto the $z$-axis: each cell is assigned
to Z-plane $z = x_3$, partitioning the $2^{Dn}$ cells into $n_z = 2^n$
planes of $2^{(D-1)n}$ cells each.

For a scalar function $f$ on the curve, let $\mu_z$ be the mean of $f$ on
plane $z$. Define the spatial covariance between planes using the $2D$
face-neighbors of each cell:

\begin{equation}
C[z_1, z_2] = \frac{1}{N_{z_1,z_2}} \sum_{\text{cells } c}
\sum_{\text{neighbors } c'} (f(d_c) - \mu_{z_1})(f(d_{c'}) - \mu_{z_2})
\label{eq:cov}
\end{equation}

where $d_c$ is the Hilbert index of cell $c$, the inner sum is over the
$2D$ face-neighbors, and $N_{z_1,z_2}$ is the number of contributing pairs.

\textbf{Observation.} $C[z_1, z_2] = 0$ for all $|z_1 - z_2| > 1$.

\textit{Reasoning.} A face-neighbor changes exactly one coordinate by
$\pm 1$. The Z-coordinate is $x_3$. For the $2(D-1)$ neighbors in
directions other than $\pm z$, the Z-coordinate is unchanged --- the
neighbor stays in the same plane. For the 2 neighbors in the $\pm z$
direction, the Z-coordinate changes by $\pm 1$ --- the neighbor moves to
an adjacent plane. No face-neighbor changes the Z-coordinate by 2 or more.
Therefore no spatial neighbor pair connects planes at distance $d > 1$.

This holds for any dimension $D \geq 3$, any order $n$, and any function
$f$. It is a geometric property of face-adjacency, not a statistical
property of $f$.

\section{Consequences}

\subsection{Tridiagonal Toeplitz Structure}

Since $C[z_1, z_2]$ depends only on $|z_1 - z_2|$ (the underlying grid is
translation-invariant in the interior), the matrix is Toeplitz:
$C[z,z] \approx c_0$ for all $z$, and $C[z,z\pm 1] \approx c_1$ for all
$z$. With all other entries zero, $C$ is tridiagonal Toeplitz.

For the prime indicator $f = 1_P$ in 4D, we measure:

\begin{center}
\begin{tabular}{cccc}
\toprule
$n$ & $c_0$ & $c_1$ & $|c_1/c_0|$ \\
\midrule
4  & $-9.978 \times 10^{-3}$ & $-9.946 \times 10^{-3}$ & 0.997 \\
5  & $-6.123 \times 10^{-3}$ & $-6.120 \times 10^{-3}$ & 0.999 \\
6  & $-4.128 \times 10^{-3}$ & $-4.128 \times 10^{-3}$ & 1.000 \\
7  & $-2.971 \times 10^{-3}$ & $-2.971 \times 10^{-3}$ & 1.000 \\
\bottomrule
\end{tabular}
\end{center}

The ratio $|c_1/c_0| \approx 1.0$ follows from the normalization: the
diagonal averages over $2(D-1) = 6$ same-plane neighbors per cell, while
the off-diagonal averages over 2 $\pm z$ neighbors; after dividing by
pair count, both estimate $\text{Var}[f]$. The ratio approaching 1.0
confirms correct normalization.

\subsection{Eigendecomposition}

The eigenvalues and eigenvectors of a tridiagonal Toeplitz matrix are
known in closed form \cite{grenander1958, kulkarni1999}:

\begin{equation}
\lambda_k = c_0 + 2c_1\cos\left(\frac{\pi k}{n+1}\right), \quad
v_k[j] = \sin\left(\frac{\pi k j}{n+1}\right), \quad k,j = 1,\ldots,n
\label{eq:eig}
\end{equation}

where $n = n_z = 2^{\text{order}}$. This is the eigendecomposition of the
discrete 1D Laplacian with Dirichlet boundaries. The Hilbert curve's
Z-plane projection is spectrally equivalent to a 1D tight-binding chain,
regardless of the ambient dimension $D$.

Table~\ref{tab:verify} compares Eq.~\ref{eq:eig} with numerical eigenvalues
for 4D spatial covariance using $f = 1_P$.

\begin{table}[h]
\centering
\begin{tabular}{ccccc}
\toprule
$n$ & $n_z$ & $r$ (analytic vs.\ numerical) & MAD (raw) & MAD (normalized) \\
\midrule
4 & 16   & 0.99986 & $1.42 \times 10^{-4}$ & 0.0036 \\
5 & 32   & 0.99992 & $5.34 \times 10^{-5}$ & 0.0021 \\
6 & 64   & 0.99996 & $2.18 \times 10^{-5}$ & 0.0013 \\
7 & 128  & 0.98499 & $1.02 \times 10^{-4}$ & 0.0084 \\
\bottomrule
\end{tabular}
\caption{Verification of Eq.~\ref{eq:eig} against numerical diagonalization.
$r$ is the Pearson correlation between the $n_z$ analytic and numerical
eigenvalues sorted by absolute magnitude. The slight degradation at $n=7$
is caused by boundary cells ($\sim$6\% of total) perturbing the strict
Toeplitz property at matrix edges.}
\label{tab:verify}
\end{table}

The eigendecomposition means that for any function $f$ on the Hilbert
curve, its Z-plane covariance can be diagonalized analytically --- no
numerical eigendecomposition is required, even for large $n$.

\subsection{Diagnostic for Curve Continuity}

A candidate space-filling curve can be tested for spatial discontinuities
by computing its Z-plane covariance for a random indicator $f$. If
$|C[z, z\pm 2]| > 0$ above machine precision, the curve has non-local
jumps between Z-planes --- it is not truly face-adjacent.

This test is cheaper than verifying adjacency for every consecutive pair
along the curve, and more sensitive than measuring the overall adjacency
percentage: a curve with 95\% adjacency but systematic boundary-crossing
jumps still produces non-zero $c_2$.

We used this test to detect an implementation error in our own 4D Hilbert
curve (Section~\ref{sec:example}), where the original ad-hoc algorithm
achieved only 50\% adjacency and produced $|c_2| \sim 10^{-3}$.

\subsection{Bounded Cross-Plane Correlation}

For spatial indexing applications \cite{moon2001}, the tridiagonal
structure bounds the ``leakage'' of range queries across Z-planes:
planes at distance $d > 1$ share \emph{no} spatial neighbor pairs. The
adjacency graph of Z-planes is exactly a 1D chain, not a higher-dimensional
graph. A query spanning planes $z_a$ through $z_b$ only has cross-talk
between adjacent planes; there is no indirect spatial coupling across
larger separations.

This is a stronger guarantee than the general locality bounds for Hilbert
curves \cite{moon2001}, which bound the number of clusters but do not
constrain the correlation between non-adjacent projections.

\section{An Application: Detecting a Broken Curve}
\label{sec:example}

During an investigation of prime distributions on Hilbert curves
\cite{wesson2026}, we implemented a 4D Hilbert curve using an ad-hoc
state machine. The curve was a bijection (every cell visited exactly once)
and superficially appeared correct. However, the Z-plane covariance matrix
had non-zero $c_2$ entries ($\sim$10$^{-3}$), flagging a problem.

The measured spatial adjacency was 50.0\% --- exactly half of consecutive
integers mapped to non-face-adjacent cells. The curve was not a true
space-filling curve; it had the same class of implementation error
(incomplete octant-boundary reflection logic) that produces non-local
jumps.

Replacing the ad-hoc algorithm with the verified Butz algorithm
\cite{butz1971, hamilton2006} restored 94.2\% adjacency and produced
$c_2 = 0$ to machine precision. The tridiagonal property served as an
immediate diagnostic that the curve was broken, without requiring a
full adjacency audit.

\section{Discussion}

The tridiagonal property is simple enough that it barely qualifies as a
result --- it follows in one sentence from the definition of
face-adjacency. The eigenvalue formula is a standard result from the 1950s
\cite{grenander1958}.

What may be new is the \emph{connection} between these two facts: that
projecting the Hilbert curve onto its Z-axis and computing spatial
covariance yields a matrix whose complete spectral decomposition is known
analytically. If this connection has been noted before, we have not found
it in the literature. If it has not, then the contribution of this note is
pointing out something that was always true but never written down.

Whether the practical consequences --- the diagnostic test and the
cross-plane correlation bound --- are useful depends on whether anyone
needs to characterize the correlation structure of Hilbert curve
projections. They arose from a number-theoretic investigation
\cite{wesson2026} where the covariance matrix was used to test for
spectral correspondence with the Riemann zeta function. In that context,
the tridiagonal structure explained why the eigenvalue spectrum took the
form it did, but did not lead to a proof of the Riemann Hypothesis.

\begin{thebibliography}{99}

\bibitem{butz1971}
A.~R.~Butz, ``Alternative Algorithm for Hilbert's Space-Filling Curve,''
\emph{IEEE Trans.\ Computers}, vol.~C-20, no.~4, pp.~424--426, 1971.

\bibitem{grenander1958}
U.~Grenander and G.~Szeg\H{o}, \emph{Toeplitz Forms and Their
Applications}, University of California Press, 1958.

\bibitem{hamilton2006}
C.~Hamilton, ``Compact Hilbert Indices,'' Dalhousie University, Technical
Report CS-2006-07, 2006.

\bibitem{hilbert1891}
D.~Hilbert, ``Ueber die stetige Abbildung einer Linie auf ein
Fl\"{a}chenst\"{u}ck,'' \emph{Mathematische Annalen}, vol.~38, pp.~459--460,
1891.

\bibitem{kulkarni1999}
D.~Kulkarni, D.~Schmidt, and S.-K.~Tsui, ``Eigenvalues of tridiagonal
pseudo-Toeplitz matrices,'' \emph{Linear Algebra and its Applications},
vol.~297, pp.~63--80, 1999.

\bibitem{moon2001}
B.~Moon, H.~V.~Jagadish, C.~Faloutsos, and J.~H.~Saltz, ``Analysis of the
Clustering Properties of the Hilbert Space-Filling Curve,'' \emph{IEEE
Trans.\ Knowledge and Data Engineering}, vol.~13, no.~1, pp.~124--141,
2001.

\bibitem{wesson2026}
R.~Wesson, ``The Hilbert Curve Prime Operator: A Complete Investigation,''
2026.

\end{thebibliography}

\end{document}
